% Môn: Vật lý
% Lớp: 10
% Chương: Chương I. Mở đầu
% Bài: L10C1 Bài 2. Các quy tắc an toàn trong phòng thực hành Vật lí · Giải nghĩa các kí hiệu kĩ thuật và biển cảnh báo an toàn trong phòng thực hành
% Số câu: 185
% App đọc file này trực tiếp — sửa rồi Commit trên GitHub, bấm Đồng bộ GitHub .tex

% L10C1 Bài 2. Các quy tắc an toàn trong phòng thực hành Vật lí



% ===== Câu 113 =====
% ID: L10C1B2-32-TN
% Mức: NB
\dangbt{Giải nghĩa các kí hiệu kĩ thuật và biển cảnh báo an toàn trong phòng thực hành}
\begin{ex}
% ID: L10C1B2-32-TN
% Mức: NB
Kí hiệu ngọn lửa trên nhãn hóa chất hoặc thiết bị thường cảnh báo điều gì?
\choice
{\True Chất hoặc khu vực dễ cháy}
{Khu vực có điện áp cao}
{Dụng cụ có cạnh sắc}
{Khu vực có bức xạ ion hóa}
\loigiai{
Kí hiệu ngọn lửa báo hiệu chất, vật liệu hoặc khu vực có nguy cơ cháy, cần tránh nguồn nhiệt, tia lửa và ngọn lửa trần.
}
\end{ex}

% ===== Câu 114 =====
% ID: L10C1B2-33-TN
% Mức: NB
\begin{ex}
% ID: L10C1B2-33-TN
% Mức: NB
Biển cảnh báo có hình tia sét trong tam giác thường dùng để chỉ nguy cơ nào trong phòng thực hành?
\choice
{Nguy cơ trượt ngã}
{Nguy cơ bỏng nhiệt}
{\True Nguy cơ điện giật}
{Nguy cơ vỡ dụng cụ thủy tinh}
\loigiai{
Biển có hình tia sét dùng để cảnh báo khu vực hoặc thiết bị có nguy cơ gây điện giật, cần tránh tiếp xúc trực tiếp và tuân thủ quy tắc an toàn điện.
}
\end{ex}

% ===== Câu 116 =====
% ID: L10C1B2-34-TN
% Mức: TH
\begin{ex}
% ID: L10C1B2-34-TN
% Mức: TH
Ý nghĩa chung của các biển cảnh báo an toàn trong phòng thực hành là gì?
\choice
{Trang trí cho phòng thực hành dễ quan sát hơn}
{Chỉ dùng để phân loại dụng cụ thí nghiệm}
{\True Nhắc nhở người học nhận biết nguy cơ và thực hiện đúng quy tắc an toàn}
{Thay thế hoàn toàn hướng dẫn của giáo viên}
\loigiai{
Biển cảnh báo giúp người học nhận biết nhanh nguy cơ mất an toàn, từ đó thao tác cẩn thận và tuân thủ hướng dẫn của giáo viên.
}
\end{ex}

